\chapter{SVAR with Long-Run Restrictions}
\label{cap:svar2}
% ── index ───────────────────────────────────────────────────────
\index{Blanchard-Quah, decomposition of|textbf}
\index{permanent shock|textbf}
\index{transitory shock|textbf}
\index{long-run restrictions|textbf}
\index{SVEC}

% ─────────────────────────────────────────────────────────────────────────────
\section{Blanchard-Quah Identification}

Ch.~\ref{cap:svar} (Cholesky SVAR) imposes \emph{short-run} restrictions:
the variable in position 1 is not contemporaneously affected by the second.
Blanchard and Quah (1989) propose an alternative via \emph{long-run}
restrictions: permanent and transitory shocks.

The central idea: in macroeconomics, a supply shock has a permanent effect
on output ($y$), while a demand shock has only a transitory effect. The restriction:
\[
  \sum_{h=0}^\infty \Theta_h^{(1,1)} = \text{free}, \quad
  \sum_{h=0}^\infty \Theta_h^{(1,2)} = 0
\]
imposes that shock 2 does not affect $y$ in the long run. This restriction identifies
the decomposition $B$ without any contemporaneous impact restriction.

\subsection*{BQ Decomposition}

Given the reduced-form VAR $y_t = \Phi_1 y_{t-1} + u_t$:
\begin{enumerate}
  \item Compute the sum $C = (I - \Phi_1)^{-1}$ — long-run multiplier
  \item Factor $C\,\Sigma_u\,C' = F\,F'$ (Cholesky of $F$)
  \item Extract $B = C^{-1} F$ — the long-run structural shocks
\end{enumerate}

% ─────────────────────────────────────────────────────────────────────────────
\section{DGP Verification}

DGP: bivariate VAR(1) where $y_1$ is a random walk (shock 1 permanent),
$y_2$ is stationary and responds transitorialy to shock 1.

\begin{lstlisting}[caption={SVAR Blanchard-Quah}]
seed(42)
// y1_t = y1_{t-1} + e1_t               (random walk)
// y2_t = 0.5*y2_{t-1} + 0.4*e1_t + 0.6*e2_t  (stationary)

let m_bq = svar(df, y1, y2, lags=1, id=longrun)
display m_bq
let m_irf_bq = sirf(m_bq, steps=20)
display m_irf_bq
\end{lstlisting}

\subsection*{B Matrix (structural shocks)}

\begin{lstlisting}[language={}, basicstyle=\ttfamily\small, frame=none,
  numbers=none, backgroundcolor=\color{codbg}]
SVAR(1) — Long-run restrictions
B Matrix:
[ 0.4907   0.0358 ]
[ 0.1935   0.3031 ]
\end{lstlisting}

The matrix $B$ is not triangular — unlike Cholesky, long-run restrictions
do not imply contemporaneous impact restrictions.

\subsection*{Structural IRF}

\begin{lstlisting}[language={}, basicstyle=\ttfamily\small, frame=none,
  numbers=none, backgroundcolor=\color{codbg}]
Impulse: Shock 1 (permanent)
  h= 0:  y1=+0.491  y2=+0.194
  h= 1:  y1=+0.464  y2=+0.087
  h= 5:  y1=+0.405  y2=-0.003   (effect on y2 decays)
  h=19:  y1=+0.271  y2=-0.005   (y1 retains permanent effect)

Impulse: Shock 2 (transitory)
  h= 0:  y1=+0.036  y2=+0.303
  h= 1:  y1=+0.016  y2=+0.144
  h= 5:  y1=-0.001  y2=+0.007   (effect on y1 decays to ~0)
  h=19:  y1=-0.001  y2= 0.000   (transitory confirmed) ✓
\end{lstlisting}

Shock 2 affects $y_1$ only transitorialy (sum $\to 0$), confirming the
long-run restriction. Shock 1 has a growing cumulative effect on $y_1$
— consistent with the unit root.

% ─────────────────────────────────────────────────────────────────────────────
\section{Syntax}

\begin{lstlisting}
svar(df, y1, y2, lags=1, id=longrun)     // BQ: long-run restrictions
svar(df, y1, y2, lags=1, id=cholesky)    // Cholesky: contemporaneous restrictions
sirf(m_svar, steps=20)                   // structural IRF
\end{lstlisting}

\begin{nota}
  BQ identification requires the variables to be integrated of order 1
  ($I(1)$). For stationary series, $(I - \hat{\Phi}_1)^{-1}$ is finite and
  the permanent/transitory decomposition has no natural interpretation.
  Verify unit roots (ch.~\ref{cap:arima}, \ref{cap:raiz_unitaria2})
  before applying BQ.
\end{nota}
